\section{Introduction}

Gaussian process \index{Gaussian process}es (GPs) provide a principled, non-parametric approach to regression and classification. They define a distribution over functions, with the covariance structure encoding assumptions about smoothness, periodicity, and other properties. Kriging \index{Kriging} (geostatistics) is the application of GPs to spatial interpolation, named after D.G. Krige and formalized by Matheron.

In this chapter, we cover:
\begin{itemize}
    \item Definition and properties of Gaussian processes
    \item Covariance functions (kernel \index{kernel (Gaussian process)}s)
    \item Gaussian process regression (Kriging)
    \item Hyperparameter optimization
    \item Scalable approximations for large datasets
    \item Applications to spatial statistics and time series
\end{itemize}

\section{Gaussian Processes}

\begin{definition}[Gaussian Process]
A Gaussian process is a collection of random variables $\{f(x) : x \in \mathcal{X}\}$ such that any finite subset $(f(x_1), \ldots, f(x_n))$ has a multivariate normal distribution.
\end{definition}

A GP is completely specified by its mean function $m(x) = \E[f(x)]$ and covariance function \index{covariance function} $k(x, x') = \Cov(f(x), f(x'))$:
\[
f \sim \mathcal{GP}(m, k).
\]

\section{Covariance Functions (Kernels)}

Common kernels:
\begin{itemize}
    \item \textbf{Squared Exponential (RBF):} $k(x,x') = \sigma^2 \exp\left(-\frac{\|x-x'\|^2}{2\ell^2}\right)$
    \item \textbf{Matérn \index{kernel!Matérn}:} $k(x,x') = \sigma^2 \frac{2^{1-\nu}}{\Gamma(\nu)} \left(\frac{\sqrt{2\nu}\|x-x'\|}{\ell}\right)^\nu K_\nu\left(\frac{\sqrt{2\nu}\|x-x'\|}{\ell}\right)$
    \item \textbf{Periodic:} $k(x,x') = \sigma^2 \exp\left(-\frac{2\sin^2(\pi|x-x'|/p)}{\ell^2}\right)$
    \item \textbf{Rational Quadratic:} $k(x,x') = \sigma^2 \left(1 + \frac{\|x-x'\|^2}{2\alpha\ell^2}\right)^{-\alpha}$
    \item \textbf{Linear:} $k(x,x') = \sigma^2 x^T x' + c$
\end{itemize}

\section{Gaussian Process Regression}

Given observations $\mathbf{y} = (y_1, \ldots, y_n)^T$ at inputs $\mathbf{X} = (x_1, \ldots, x_n)^T$, with noise $y_i = f(x_i) + \epsilon_i$, $\epsilon_i \sim \mathcal{N}(0, \sigma_n^2)$:
\[
\mathbf{y} \sim \mathcal{N}(m(\mathbf{X}), K(\mathbf{X}, \mathbf{X}) + \sigma_n^2 I).
\]

For test inputs $\mathbf{X}_*$, the predictive distribution is
\[
f_* | \mathbf{y}, \mathbf{X}, \mathbf{X}_* \sim \mathcal{N}(\mu_*, \Sigma_*),
\]
where
\begin{align*}
\mu_* &= m(\mathbf{X}_*) + K(\mathbf{X}_*, \mathbf{X}) [K(\mathbf{X}, \mathbf{X}) + \sigma_n^2 I]^{-1} (\mathbf{y} - m(\mathbf{X})) \\
\Sigma_* &= K(\mathbf{X}_*, \mathbf{X}_*) - K(\mathbf{X}_*, \mathbf{X}) [K(\mathbf{X}, \mathbf{X}) + \sigma_n^2 I]^{-1} K(\mathbf{X}, \mathbf{X}_*).
\end{align*}

\section{Log Marginal Likelihood}

\[
\log p(\mathbf{y} | \mathbf{X}, \theta) = -\frac{1}{2} \mathbf{y}^T K_\theta^{-1} \mathbf{y} - \frac{1}{2} \log |K_\theta| - \frac{n}{2} \log 2\pi,
\]
where $K_\theta = K(\mathbf{X}, \mathbf{X}) + \sigma_n^2 I$. Hyperparameters $\theta$ are optimized by gradient ascent.

\section{Python Implementation}

In \texttt{python/gp.py}:

\begin{lstlisting}[style=python, caption={Key functions in python/gp.py}]
class GaussianProcess:
    def __init__(self, kernel, noise=1e-6)
    def fit(self, X, y)
    def predict(self, X_star, return_std=True)
    def log_marginal_likelihood(self, theta)
    def optimize_hyperparameters(self, X, y)

rbf_kernel(x1, x2, length_scale, variance)
matern_kernel(x1, x2, length_scale, variance, nu)
periodic_kernel(x1, x2, length_scale, variance, period)
kriging(x_obs, y_obs, x_pred, kernel)  # Simple Kriging
ordinary_kriging(x_obs, y_obs, x_pred) # Ordinary Kriging
\end{lstlisting}

\section{Numerical Examples}

\begin{example}[GP Regression]
\begin{lstlisting}[style=python]
>>> from gp import GaussianProcess, rbf_kernel
>>> import numpy as np
>>> X = np.array([0, 1, 2, 3, 4]).reshape(-1, 1)
>>> y = np.sin(X.ravel()) + 0.1 * np.random.randn(5)
>>> gp = GaussianProcess(rbf_kernel, noise=0.01)
>>> gp.fit(X, y)
>>> X_test = np.linspace(0, 4, 50).reshape(-1, 1)
>>> mu, std = gp.predict(X_test, return_std=True)
\end{lstlisting}
\end{example}

\section{Exercises}

\begin{exercise}
Implement the Matérn kernel with $\nu = 3/2$ and $\nu = 5/2$ using the simplified form.
\end{exercise}

\begin{exercise}
Derive the gradient of the log marginal likelihood \index{marginal likelihood} with respect to kernel hyperparameters.
\end{exercise}

\begin{exercise}
Implement sparse Gaussian process approximations (FITC, VFE) for large $n$.
\end{exercise}

\begin{exercise}
Apply ordinary Kriging to a spatial dataset and compare with inverse distance weighting.
\end{exercise}



\section{Exercise Solutions}

\begin{solution}
A Gaussian process $f \sim \mathcal{GP}(m(x), k(x,x'))$ is determined by its mean function $m(x)$ and covariance function $k(x,x')$. Any finite collection $(f(x_1), \ldots, f(x_n))$ is jointly Gaussian with mean $(m(x_1), \ldots, m(x_n))$ and covariance matrix $K_{ij} = k(x_i, x_j)$.
\end{solution}

\begin{solution}
The RBF kernel $k(x,x') = \sigma^2 \exp(-\|x-x'\|^2/(2\ell^2))$ produces infinitely differentiable functions. The length scale $\ell$ controls smoothness: small $\ell$ gives rapid oscillations, large $\ell$ gives smooth functions. The signal variance $\sigma^2$ controls the amplitude.
\end{solution}


\section{References}

\begin{itemize}
    \item \cite{rasmussen-williams} -- Rasmussen \& Williams, \textit{Gaussian Processes for Machine Learning}
    \item \cite{cressie} -- Cressie, \textit{Statistics for Spatial Data}
    \item \cite{gauss-theoria} -- Gauss's work on interpolation and geodesy
    \item \cite{stein} -- Stein, \textit{Interpolation of Spatial Data}
\end{itemize}